\input{../tex-inputs/latex-leseansicht-vorspann}

\section[Arthur Schnitzler an Theodor von Sosnosky, 26. 5. 1901]{L04234 Arthur Schnitzler an Theodor von Sosnosky, 26. 5. 1901}
\nopagebreak\mylabel{L04234v}
\rehead{ }\normalsize\beginnumbering\briefempfaengerindex{Sosnosky, Theodor von@\textsc{Sosnosky, Theodor von}!zzzSchnitzler, Arthur@\emph{von Arthur Schnitzler}!1901-05-261@{26. 5. 1901}|(be}
\toendnotes[C]{\smallbreak\pagebreak[2]}
\correspDesc{Versand  durch Arthur Schnitzler am 26. 5. 1901 in Wien
\newline{}Erhalt  durch Theodor von Sosnosky am 26. 5. 1901 \textbf{Ort fehlend} }\toendnotes[C]{\smallbreak}
\Standort{CUL, Schnitzler, B 126.}
\physDesc{Brief, 1 Blatt, 4 Seiten
\newline{}Handschrift: , deutsche Kurrent
\newline{}Zusatz: Markierungen einer Versteigerung: »5468« und »15–« }
\buchAbdrucke{\weitereDrucke{1) Theodor v. Sosnosky: \emph{Unveröffentlichte Schnitzler-Briefe über die »Leutnant-Gustl«-Affäre. Eine Sensation vor dreißig Jahren.} In: \emph{Neues Wiener Journal}, Jg. 39, Nr. 13.624, 26. 10. 1931, S. 4.} \weitereDrucke{2) Ursula Renner: \emph{Dokumentation eines Skandals Arthur Schnitzlers »Lieutenant Gustl«.} In: \emph{Hofmannsthal-Jahrbuch}, Jg. 15 (2007), S. 210–211.} }\toendnotes[C]{\smallbreak}
\pstart
           \raggedleft{}{\pb}Wien\oindex{Wien@\textbf{Wien}, \emph{Verwaltungsgebiet}|pw}, 26. 5. 1901.\pend
           \vspace{0.5em}
\pstart
           Verehrteſter Herr von Sosnosky, ich danke Ihnen herzlich für die
               freundlichen Mittheilungen und hoffe \label{K_L04234-1v}\edtext{bald leſen zu dürfen}{\lemma{\textnormal{\emph{bald lesen zu dürfen}}}\Cendnote{\textnormal{Die Veröffentlichung
                  der Rezension verzögerte sich um mehrere Monate: Th. v. S.\pwindex{Sosnosky, Theodor von 4.\,1.\,1866 Budapest – 3.\,2.\,1943 Wien@\textsc{Sosnosky, Theodor von} (4.\,1.\,1866 Budapest – 3.\,2.\,1943 Wien), \emph{Schriftsteller}|pwk} [= Theodor von Sosnosky\pwindex{Sosnosky, Theodor von 4.\,1.\,1866 Budapest – 3.\,2.\,1943 Wien@\textsc{Sosnosky, Theodor von} (4.\,1.\,1866 Budapest – 3.\,2.\,1943 Wien), \emph{Schriftsteller}|pwk}]: \emph{Bücher und Zeitschriftenschau}\pwindex{Sosnosky, Theodor von 4.\,1.\,1866 Budapest – 3.\,2.\,1943 Wien@\textsc{Sosnosky, Theodor von} (4.\,1.\,1866 Budapest – 3.\,2.\,1943 Wien), \emph{Schriftsteller}!Bücher und Zeitschriftenschau [Frau Bertha Garlan]@\strich\emph{Bücher und Zeitschriftenschau [Frau Bertha Garlan]}|pwk}. In: \emph{Norddeutsche Allgemeine Zeitung}\pwindex{Norddeutsche Allgemeine Zeitung@\emph{Norddeutsche Allgemeine Zeitung}|pwk}, Nr. 50, 28. 2. 1902, Beilage.}}}\label{K_L04234-1}, was Sie über die »Garlan\pwindex{Schnitzler, Arthur 15.\,5.\,1862 Wien – 21.\,10.\,1931 ebd.@\textsc{Schnitzler, Arthur} (15.\,5.\,1862 Wien – 21.\,10.\,1931 ebd.), \emph{Schriftsteller, Mediziner}!Frau Bertha Garlan. Roman@\strich\emph{Frau Bertha Garlan. Roman}|pw}« geſagt haben. Ich{ }ſchließe gewiſs nicht unrichtig, dſs es das Neue Wiener Tagblatt\orgindex{Neues Wiener Tagblatt@Neues Wiener Tagblatt|pw} war, welches Ihre Kritik als zu freundlich
               refüſirt hat. Ihren Bemerkungen zum Lieutn Guſtl\pwindex{Schnitzler, Arthur 15.\,5.\,1862 Wien – 21.\,10.\,1931 ebd.@\textsc{Schnitzler, Arthur} (15.\,5.\,1862 Wien – 21.\,10.\,1931 ebd.), \emph{Schriftsteller, Mediziner}!Lieutenant Gustl. Novelle@\strich\emph{Lieutenant Gustl. Novelle}|pw}
               müſſen Sie mir geſtatten zwei Worte zu erwidern. Ich finde, daſs man mir mit
               demſelben Recht als \introOben{}man\introOben{} mir Feind{\pb}ſeligkeit gegen den Offiziersſtand vorwirft, weil
               ich den Guſtl\pwindex{Schnitzler, Arthur 15.\,5.\,1862 Wien – 21.\,10.\,1931 ebd.@\textsc{Schnitzler, Arthur} (15.\,5.\,1862 Wien – 21.\,10.\,1931 ebd.), \emph{Schriftsteller, Mediziner}!Lieutenant Gustl. Novelle@\strich\emph{Lieutenant Gustl. Novelle}|pw} geſchrieben – Gehäſſigkeit gegen
               die Aerzte vorwerfen könnte, weil ich den Ferdinand Schmidt\pwindex{Schnitzler, Arthur 15.\,5.\,1862 Wien – 21.\,10.\,1931 ebd.@\textsc{Schnitzler, Arthur} (15.\,5.\,1862 Wien – 21.\,10.\,1931 ebd.), \emph{Schriftsteller, Mediziner}!Vermächtnis. Schauspiel in drei Akten@\strich\emph{Das Vermächtnis. Schauspiel in drei Akten}|pwv} im Vermächtnis\pwindex{Schnitzler, Arthur 15.\,5.\,1862 Wien – 21.\,10.\,1931 ebd.@\textsc{Schnitzler, Arthur} (15.\,5.\,1862 Wien – 21.\,10.\,1931 ebd.), \emph{Schriftsteller, Mediziner}!Vermächtnis. Schauspiel in drei Akten@\strich\emph{Das Vermächtnis. Schauspiel in drei Akten}|pw}, Gehäſſigkeit gegen den Liberalismus, weil ich den Loſatti\pwindex{Schnitzler, Arthur 15.\,5.\,1862 Wien – 21.\,10.\,1931 ebd.@\textsc{Schnitzler, Arthur} (15.\,5.\,1862 Wien – 21.\,10.\,1931 ebd.), \emph{Schriftsteller, Mediziner}!Vermächtnis. Schauspiel in drei Akten@\strich\emph{Das Vermächtnis. Schauspiel in drei Akten}|pwv} (im gleichen Stück\pwindex{Schnitzler, Arthur 15.\,5.\,1862 Wien – 21.\,10.\,1931 ebd.@\textsc{Schnitzler, Arthur} (15.\,5.\,1862 Wien – 21.\,10.\,1931 ebd.), \emph{Schriftsteller, Mediziner}!Vermächtnis. Schauspiel in drei Akten@\strich\emph{Das Vermächtnis. Schauspiel in drei Akten}|pwv}), Gehäſſigkeit gegen
               Theaterdirektoren (Schneider\pwindex{Schnitzler, Arthur 15.\,5.\,1862 Wien – 21.\,10.\,1931 ebd.@\textsc{Schnitzler, Arthur} (15.\,5.\,1862 Wien – 21.\,10.\,1931 ebd.), \emph{Schriftsteller, Mediziner}!Freiwild. Schauspiel in 3 Akten@\strich\emph{Freiwild. Schauspiel in 3 Akten}|pwv}
               in Freiwild\pwindex{Schnitzler, Arthur 15.\,5.\,1862 Wien – 21.\,10.\,1931 ebd.@\textsc{Schnitzler, Arthur} (15.\,5.\,1862 Wien – 21.\,10.\,1931 ebd.), \emph{Schriftsteller, Mediziner}!Freiwild. Schauspiel in 3 Akten@\strich\emph{Freiwild. Schauspiel in 3 Akten}|pw}) gegen Weinhändler (Schwager Garlan\pwindex{Schnitzler, Arthur 15.\,5.\,1862 Wien – 21.\,10.\,1931 ebd.@\textsc{Schnitzler, Arthur} (15.\,5.\,1862 Wien – 21.\,10.\,1931 ebd.), \emph{Schriftsteller, Mediziner}!Frau Bertha Garlan. Roman@\strich\emph{Frau Bertha Garlan. Roman}|pw}) u. ſ. w. vorwerfen konnte. In jedem Stand
               gibt es {\pb}eine große Anzahl minderwerthiger Individuen und ich erkenne keinem Stande
               das Recht zu,{ }ſich dieſe Conſtatirung zu verbieten. Und dabei möchte ich für meinen
               Fall noch erwähnen, dſs mein »Guſtl\pwindex{Schnitzler, Arthur 15.\,5.\,1862 Wien – 21.\,10.\,1931 ebd.@\textsc{Schnitzler, Arthur} (15.\,5.\,1862 Wien – 21.\,10.\,1931 ebd.), \emph{Schriftsteller, Mediziner}!Lieutenant Gustl. Novelle@\strich\emph{Lieutenant Gustl. Novelle}|pw}« ein ganz
               netter, nur durch Standesvorurtheile verwirrter Burſch ist, der mit den Jahren gewiſs
               ein tüchtiger und anſtändiger Offizier werden dürfte. Aber – hievon abgeſehen, müßte
               ich mir als Dichter jedenfalls das Recht {\pb}\label{K_L04234-2v}\edtext{vindiciren}{\lemma{\textnormal{\emph{vindiciren}}}\Cendnote{\textnormal{veraltet: sich das Recht herausnehmen}}}\label{K_L04234-2}, einen \substVorne{}\textsuperscript{Dichter}\substDazwischen{}Offizier\substHinten{}{ }\introOben{}ſelbst\introOben{} zu erfinden, der gewohnheitsmäßig{ }ſeine Großcouſinen umbringt; we{\geminationn}
               ich einen{ }ſolchen Offizier für meine Novelle brauchte.– Ihnen, mein verehrter Herr
               von Sosnosky, dem ich{ }ſchon für{ }ſoviele wahre Liebenswürdigkeit und für{ }ſoviel
               Verständnis zu danken habe, darf ich heute auch dieſe meine Stellung zu der von Ihnen
               angedeuteten Frage andeuten, ohne die Befürchtg misverſtanden zu werden.\pend
           
\pstart
           Mit den verbindlichſten Grüßen {\\[\baselineskip]}Ihr sehr ergebner \spacefill\mbox{Arthur
                  Schnitzler.}\pend
           \leftskip=0em{}\selectlanguage{ngerman}\endnumbering\briefempfaengerindex{Sosnosky, Theodor von@\textsc{Sosnosky, Theodor von}!zzzSchnitzler, Arthur@\emph{von Arthur Schnitzler}!1901-05-261@{26. 5. 1901}|)be}\mylabel{L04234h}
\begin{anhang}
\end{anhang}\newcommand{\dateiname}{L04234}\newcommand{\titel}{Arthur Schnitzler an Theodor von Sosnosky, 26. 5. 1901}\newcommand{\editorInnen}{Herausgegeben von Jahnke, SelmaMüller, Martin Anton}\input{../tex-inputs/latex-leseansicht-abspann}
